\documentclass[12pt, a4paper]{article}
\usepackage[
    top=3 cm,
    left=3cm,
    bottom=2cm,
    right=2cm
]{geometry}

\usepackage[utf8]{inputenc}
\usepackage[T1]{fontenc}
\usepackage[brazilian]{babel}
\usepackage{lipsum}
\usepackage{setspace}
\usepackage{enumitem}
\usepackage{setspace}
\usepackage{graphicx}
\usepackage{booktabs}
\usepackage{float}
\usepackage{amsmath}
\usepackage{amsfonts}
\usepackage{amssymb}

\onehalfspacing
%\singlespacing
%\doublespacing



\setlength{\parindent}{1.25cm}



\title{Curso Final\LaTeX}
\author{GUILHERME CARDOSO\thanks{Universidade do Estado de Santa Catarina} }
\date{\today}


\begin{document}
\maketitle
\begin{abstract}
    
\end{abstract}
\newpage
\section{Introdução}
Este documento destaca as atividades desenvolvidas durante o curso de LaTeX ministrada pelo professor Lucas Alexandre Zick, onde no decorrer das 12 aulas se teve o foco para o entendimento da ferramenta e sua familiarização para com um nivel avançado e preparação para o mundo academico e seus desafios como um todo.

A construção se deve a intuito de tornar pratico os metodos expostos trazendo assim muito mais familiaridade com os comandoss basicos, partindo desde a criação do preambulo até quais tipos de pacotes se deve impotar, dependendo da finalidade e para qual tipo de documento a utilização do LaTex.

\section{Fundamentação Teórica}


\textsf{\huge TFG}
TFG
\texttt{TFG}
\textsc{\tiny TFG}
{\scriptsize TFG}

\begin{center}
    asssssssssssssssssssssssssssssssssssssssdasddddddddd
\end{center}

\begin{flushleft}
aasdddddddddddddddddasdddasdadadadadadadadadada
\end{flushleft}

\begin{flushright}
    asdaaasdasdasdasdadadaddasasasasa
\end{flushright}



Texto1.
\vspace{2cm}
Texto2
Salve rapziada\footnote{Não salvo maix}bom dia\footnote{boa noite}


Assinatura: \hspace{9cm} Data:








\subsection{}

\subsection{Latex}


Em meados do anos 70, se teve um aumento gradativo da computação onde consequentemente as digitalizações de artigos e livros cientificos, e o resultado foi uma qualidade não agradavel para fórmulas matemáticas e o espaçamento dos textos, pu seja, foi apontado que a produação cientifica exige um rigor que os editores de texto da época não se adequavam.


Se deve apontar dois nomes para a criação do latex para dois individeios, o primeiro deles sendo o professor Donald Knuth (Universidade de Stanford), que em 1978, se viu frustado com impressões de um de seus livros, mais especificamente o The Art of Computer Programming, e assim nasceu o TeX, um sistema que busca focar em perfeição matemática e tipográfica, onde o objetivo descrito pelo proprio Donald Knuth é: "O TeX foi projetado com dois objetivos principais em mente: permitir que qualquer pessoa produza livros de alta qualidade com esforço mínimo, e fornecer um sistema que daria
exatamente os mesmos resultados em todos os computadores"\cite{}

O segundo nome a ser destacado apareceu em 1985, indetificando o potencial do TeX original, levando em conta que é considerado um dos sofwares mais estáveis ja desenvolvidos da historias, o cientista Leslie Lamport constriu uma "camada" por cima do TeX, trazendo assim comandos amigaveis e um grau de dificuldade aceitavel para usuários comuns e assim nascia o LaTeX(Lamport TeX)



































\subsection{A plataforma Overleaf}

Um dos grandes motivos para a resistência ao uso do LaTeX é a necessidade de diversas instalações e distribuições locais (como \textit{TeX Live} ou \textit{MiKTeX}), que costumam exigir de 4GB a 6GB de armazenamento. Além disso, há a complexidade de configurar editores locais (como o \textit{TeXmaker} ou o \textit{VS Code}), dependendo do nível de experiência do usuário. 

A plataforma \textit{Overleaf} destaca-se por ser uma aplicação totalmente baseada na \textit{web}. Desse modo, o usuário necessita apenas de uma conexão com a internet e de um navegador moderno para criar seus documentos. Por ter seu enfoque no processamento e no armazenamento em nuvem, a ferramenta garante fácil acesso e praticidade para múltiplos dispositivos de diferentes sistemas operacionais.

Além da portabilidade, o \textit{Overleaf}. A plataforma, permite que múltiplos autores trabalhem no mesmo arquivo simultaneamente em tempo real, facilitando o processo de revisão entre alunos e orientadores. Por fim, a plataforma disponibiliza uma vasta biblioteca de modelos (\textit{templates}) de periódicos científicos e universidades, reduzindo drasticamente a curva de aprendizado inicial e o tempo gasto com a configuração do preâmbulo do documento.






















\subsection{Leitura}

Se deve destacar as oportunidades que o micro-gerenciamento textual é permitido no LaTex, onde é possivel, destacar palavras, configurar os recuos de paragrafos, alinhar textos, adicionar notas de rodapé e criar listas(tópicos), trazendo asim possibilidades para o nosso controle no quesito de leitura.



Diferentemente dos editores tradicionais como Microsoft world nos quais a formatação é aplicada manualmente e está sujeita a dificuldade e inconsistensias visualmente falando, o LaTEX permiete a atuomatização do acabamento do texto, trazendo por exemplo uma segurança maior qunado se trata de seguir padroes para normas da ABNT por exemplo, onde serão estipuladas no preambulo por meio de pacotes que definem desde os espaçamento entre linhas, recuo de paragrafos, entre outras opções para padronizações.


Cabe o apontamento a da padronizações permissivas e objetivas para notas de rodapes e a modularidade que existe para o sumário e referencias no texto sem uma atenção de retrabalho como para ferrementas basicas como world.



\subsection{ Equações Matemáticas}





O LaTEX tem como uma das suas dominancias dentro do mundo das engenharias e ciencias exatas devido ao fato de sua particiade para com a escrita de formulas complexas, tendo enorme vantagens para seus concorrentes (como world) principalmente pelo esforço de sua visibildade comparada ao LaTEX que se resume inicialmente a utilização de tres pocotes centrais  fornecidos pela American Mathematical Society (AMS)


Os pacotes inicais sendo os Devemos adicionar isto ao nosso preâmbulo:
\begin{itemize}
    \item \textbackslash usepackage{amsmath}: Fornece os ambientes essenciais (equações alinhadas matrizes);
    \item \textbackslash usepackage{amsfonts}: Fornece fontes matemáticas 
    clássicas;
    \item \textbackslash usepackage{amssymb}: Adiciona dezenas de símbolos essenciais (≈, ≤, ≥, R, ∇)

\end{itemize}

























\section{Metodologia Prática}
\subsection{Aula01}
Foi apresentado uma introdução historica sobre a criação, utilização e o proposito inicial do \textbf{TEX} sendo o cerne do problema o fator da falta de praticidade para artigos impressos em papel, vale destacar que a sua criação foi concebida com quase 100\% livre de problemas ou bugs. Devido sua complexidade, anos depois se teve a sua adaptação para \textbf{LaTEX}, focado na sua adptação para uma utilização mais amigavel.


A utilização do \textbf{LaTEX} se destaca pela facil criação de padronizações e sua praticidade, quando comparada a ferramentas como word que são catalogadas como \textit{What You See What You Get}(WYSWYG), ou seja são visuais e se tem um cuidado maior no ambito da pesquisa principalmente pelo fator de diferentes versões que podem trazer problemas e dores de cabeça.


Partindo para ferramentas cabe destacar o OverLeaf, com sua praticidade no quesito de não ser necessaria uma instalação em si(zero setup), ou seja, basta criar uma conta e usar a ferramenta on-line pelo navegador, trazendo com isso uma variedade de acesso seja com celular ou computador, e sua capacidade de compartilhar o acesso com diferentes usuários.




\subsection{Aula02}

Se teve uma breve revisão da aula anterior com um enfoque para o \textbf{Preâmbulo} e o Corpo Textual, apartir disso se teve como objetivo o Chassi do documento ou seja o esqueleto e uma base para a escrita de um artigo ou TCC, tendo assim uma base inicial para as importações iniciais e os fundamentos para o Preambulo.


O destaque inicial se teve com os comandos inciais para o inicio do Preambulo e os comandos essenciais e quase que obrigatorios, destaca-se para as classses dos documentos sendo os possiveis Article, report e book, onde cada um tem suas vantagens e especifidades dependendo do objetivo final.


Um levantamento para destacar algumas opções de paramentros modificadores dependendo da finalidade e antes do nome da classe:
\begin{itemize}
    \item \textbackslash documentclass [12pt , a4paper , twocolumn ]\{ article \}
    \item \textbackslash documentclass [10pt , letterpaper ]\{ report \}
    \item  \textbackslash documentclass [11pt , draft ]\{ article \}

\end{itemize}

Cabe destacar a praticidade em ajsutar as margens usando o pacote \textit{gemetry} sendo destacado como o padrão absoluto para essa finalidade,se teve atenção para a existencia do Pacote \textit{Lipsum} para testes de design antes da escrita da documentação real,  foi apontado os pacotes para documentos em português:

\begin{itemize}
    \item \textbackslash usepackage [ brazil ]\{ babel \}
    \item \textbackslash usepackage [ utf 8]\{ inputenc \}
    \item  \textbackslash usepackage [T 1]\{ fontenc \}

\end{itemize}















\subsection{Aula03}
Se teve destaque para as diferentes formas dos comando para conceitos pontuais para a escrita de um documento, levantando o pontos para a utilziação e maniupalção de como destacar palavras com a finalidade desejada:

\begin{itemize}
    \item \textbackslash textbf\{Negrito\}
    \item \textbackslash textit\{italico\}
    \item \textbackslash underline\{Sublinhado\}
    \item \textbackslash textsc\{Small Caps\}
\end{itemize}













\subsection{Aula04}
\includegraphics[width=10cm, height= -10cm]{}

\begin{figure}
    \centering
    \includegraphics[width
    =0.5\linewidth]{}
    \caption{}
    \label{fig:gauches}
\end{figure}


\begin{tabular} {l|c|r}
    \hline
    \hline
    componente && salve && teste \\
    \hline
    basdadsasdasdasd  && salve && test \\
    \hline
    componente && salve && test \\
    \hline
    \hline
\end{tabular}





\subsection{Aula05}

Editores comuns(Word) acacabam sofrendo de um problema cronico perante a necessidade de inserção e imagens, onde o tamanho é o principal agravante, gerando catalaizadores para a imagem não caber no espaço e pular para página seguinte, e trazendo assim um buraco enorme onde a imagem deveria estar porem não coube, e existe um ultimo empencilho que se adicionar uma imgaem em algum paragrafo antes, todas as figuras correm o risco de sair de estarem seriamente comprometidas tanto em ordem quanto visualmente.


O laTEX destaca uma solução, que é tanto as imagens quanto as tabelas são alocadas em "contêineres flutuantes"(figure e table), ou sejam se uma não cabe em determinada página, o LaTeX vai continuar a "puxar" o texto que vai vir depois, fazendo assim adptação para que a imagem se mantenha estavel sem comprometer com buracos em branco.


Essa manipulação se obtem por conta do controle da flutuação os especificadores, onde exite os parametros para serem apontados para o compilador:


\begin{itemize}
    \item h (here): Tenta colocar exatamente onde está no código.

    \item t (top): Coloca no topo da página.

    \item b (bottom): Coloca na base da página
    \item p (page): Cria uma página exclusiva para figuras/tabelas
    \item ! (override): Força o LaTeX a ignorar restrições estéticas internas.

\end{itemize}







\subsection{Aula06}



















\subsection{Aula07}

\includegraphics[width=-5cm]{imagens/images.jpg}

\section{Considerações Finais}


\end{document}
